\documentclass[aspectratio=169]{beamer}
\usetheme{Madrid}
\usecolortheme{beaver}
\usepackage[utf8]{inputenc}
\usepackage[T1]{fontenc}
\setbeamertemplate{navigation symbols}{}
\setbeamercolor{title}{fg=white,bg=cyan!60!black}
\setbeamercolor{frametitle}{fg=white,bg=cyan!50!black}
\setbeamercolor{block title}{bg=cyan!15,fg=black}
\setbeamercolor{block body}{bg=cyan!5,fg=black}

\title{G4 — My Blocks}
\subtitle{jump(), shoot(), spawnEnemy() — powers you define}
\author{ByteForge --- Grade 7 --- Game Maker}
\date{}

\begin{document}

\begin{frame}
    \titlepage
\end{frame}

\begin{frame}{Today's Plan}
    \begin{itemize}
        \item spawnEnemy()
        \item shoot()
        \item powerUp()
        \item[$\bigstar$] BOSS --- Platformer Vertical Slice
    \end{itemize}
\end{frame}

\begin{frame}{Functions — reusable named blocks}
    \Large
    A function is defined ONCE with parameters, then CALLED many times with different values — one definition, many different outcomes.
\end{frame}


\begin{frame}{A parameter that changes behaviour}
    \Large
    A true/false parameter can make the FUNCTION ITSELF decide what to draw, using an IF inside its own definition.
\end{frame}


\begin{frame}{A text parameter, not just numbers}
    \Large
    Parameters can carry any type of value — numbers, true/false, or text — the function's job is always to use whatever was passed in.
\end{frame}


\begin{frame}{BOSS — planning in pseudocode BEFORE building}
    \Large
    Writing the logic in plain English first is called pseudocode — real programmers do this constantly for anything fiddly, like jump physics.
\end{frame}

\begin{frame}{Real-World Example}
    \Large
    A recipe you cook many times (make pancakes) is a "function" — the same steps, but the exact ingredients (parameters) change each time depending on how many people you're feeding.
\end{frame}

\begin{frame}{Watch}
    \Large
    Suggested video topic: \emph{What is a function — parameters explained for beginners}

    \vspace{1em}
    \normalsize
    Search this on YouTube together, or swap in a specific video you've
    already vetted --- do not project an unverified link cold in class.
\end{frame}

\begin{frame}{Let Them Discover}
    \Large
    Think of something you do repeatedly with small variations (making a sandwich, getting ready for school). Write the general steps once, then note which parts change each time (the "parameters").
\end{frame}

\begin{frame}{Your Turn}
    \Large
    Open the student map, find \textbf{G4 — My Blocks}, and work through it
    hands-on --- everything you just saw, they now do themselves.
\end{frame}

\end{document}
